% Options for packages loaded elsewhere
\PassOptionsToPackage{unicode}{hyperref}
\PassOptionsToPackage{hyphens}{url}
\PassOptionsToPackage{dvipsnames,svgnames,x11names}{xcolor}
%
\documentclass[
  10pt,
  letterpaper,
  onecolumn]{article}
\usepackage{amsmath,amssymb}
\usepackage{iftex}
\ifPDFTeX
  \usepackage[T1]{fontenc}
  \usepackage[utf8]{inputenc}
  \usepackage{textcomp} % provide euro and other symbols
\else % if luatex or xetex
  % fontspec, NOT unicode-math: unicode-math makes math glyphs math-active,
  % which breaks \newunicodechar mappings in text mode.
  \usepackage{fontspec}
  \defaultfontfeatures{Scale=MatchLowercase}
  \defaultfontfeatures[\rmfamily]{Ligatures=TeX,Scale=1}
\fi
\IfFileExists{lmodern.sty}{\ifPDFTeX\usepackage{lmodern}\fi}{}
\ifPDFTeX\else
  % xetex/luatex font selection
  \setmainfont[]{Latin Modern Roman}
\fi
% Use upquote if available, for straight quotes in verbatim environments
\IfFileExists{upquote.sty}{\usepackage{upquote}}{}
\IfFileExists{microtype.sty}{% use microtype if available
  \usepackage[]{microtype}
  \UseMicrotypeSet[protrusion]{basicmath} % disable protrusion for tt fonts
}{}
\makeatletter
\@ifundefined{KOMAClassName}{% if non-KOMA class
  \IfFileExists{parskip.sty}{%
    \usepackage{parskip}
  }{% else
    \setlength{\parindent}{0pt}
    \setlength{\parskip}{6pt plus 2pt minus 1pt}}
}{% if KOMA class
  \KOMAoptions{parskip=half}}
\makeatother
\usepackage{xcolor}
\usepackage[margin=0.85in]{geometry}
\setlength{\emergencystretch}{3em} % prevent overfull lines
\providecommand{\tightlist}{%
  \setlength{\itemsep}{0pt}\setlength{\parskip}{0pt}}
\setcounter{secnumdepth}{5}
\newcommand{\notestitle}{Problems --- Implementation Notes}
% ---------------------------------------------------------------------------
%  Study-notes preamble (inlined by pandoc -H, so each .tex is standalone).
%  Engine: XeLaTeX.  Class: article, 10pt, one column.
% ---------------------------------------------------------------------------
\usepackage{amsmath,amssymb}
\usepackage{newunicodechar}
\usepackage{enumitem}
\usepackage{fancyhdr}
\usepackage{titlesec}
\usepackage{microtype}

% --- fonts ------------------------------------------------------------------
% Latin Modern for text/math; DejaVu Sans Mono for code (wide Unicode coverage).
\setmonofont{DejaVu Sans Mono}[Scale=MatchLowercase]

% --- Unicode the text font cannot draw -> real math glyphs ------------------
\newunicodechar{−}{\ensuremath{-}}
\newunicodechar{·}{\ensuremath{\cdot}}
\newunicodechar{×}{\ensuremath{\times}}
\newunicodechar{÷}{\ensuremath{\div}}
\newunicodechar{±}{\ensuremath{\pm}}
\newunicodechar{≤}{\ensuremath{\leq}}
\newunicodechar{≥}{\ensuremath{\geq}}
\newunicodechar{≈}{\ensuremath{\approx}}
\newunicodechar{≠}{\ensuremath{\neq}}
\newunicodechar{≅}{\ensuremath{\cong}}
\newunicodechar{≇}{\ensuremath{\ncong}}
\newunicodechar{∞}{\ensuremath{\infty}}
\newunicodechar{√}{\ensuremath{\sqrt{\;}}}
\newunicodechar{Σ}{\ensuremath{\Sigma}}
\newunicodechar{∫}{\ensuremath{\int}}
\newunicodechar{∈}{\ensuremath{\in}}
\newunicodechar{∘}{\ensuremath{\circ}}
\newunicodechar{‖}{\ensuremath{\|}}
\newunicodechar{ℓ}{\ensuremath{\ell}}
\newunicodechar{→}{\ensuremath{\rightarrow}}
\newunicodechar{←}{\ensuremath{\leftarrow}}
\newunicodechar{↔}{\ensuremath{\leftrightarrow}}
\newunicodechar{⇒}{\ensuremath{\Rightarrow}}
\newunicodechar{⇐}{\ensuremath{\Leftarrow}}
\newunicodechar{⇔}{\ensuremath{\Leftrightarrow}}
\newunicodechar{⌈}{\ensuremath{\lceil}}
\newunicodechar{⌉}{\ensuremath{\rceil}}
\newunicodechar{⌊}{\ensuremath{\lfloor}}
\newunicodechar{⌋}{\ensuremath{\rfloor}}
\newunicodechar{°}{\ensuremath{{}^{\circ}}}
\newunicodechar{✓}{\ensuremath{\checkmark}}
% greek
\newunicodechar{α}{\ensuremath{\alpha}}
\newunicodechar{β}{\ensuremath{\beta}}
\newunicodechar{γ}{\ensuremath{\gamma}}
\newunicodechar{ε}{\ensuremath{\epsilon}}
\newunicodechar{η}{\ensuremath{\eta}}
\newunicodechar{θ}{\ensuremath{\theta}}
\newunicodechar{λ}{\ensuremath{\lambda}}
\newunicodechar{μ}{\ensuremath{\mu}}
\newunicodechar{π}{\ensuremath{\pi}}
\newunicodechar{ρ}{\ensuremath{\rho}}
\newunicodechar{σ}{\ensuremath{\sigma}}
\newunicodechar{φ}{\ensuremath{\phi}}
\newunicodechar{ω}{\ensuremath{\omega}}
% super/subscripts
\newunicodechar{¹}{\ensuremath{{}^{1}}}
\newunicodechar{²}{\ensuremath{{}^{2}}}
\newunicodechar{³}{\ensuremath{{}^{3}}}
\newunicodechar{⁴}{\ensuremath{{}^{4}}}
\newunicodechar{⁵}{\ensuremath{{}^{5}}}
\newunicodechar{⁶}{\ensuremath{{}^{6}}}
\newunicodechar{⁷}{\ensuremath{{}^{7}}}
\newunicodechar{⁸}{\ensuremath{{}^{8}}}
\newunicodechar{ⁿ}{\ensuremath{{}^{n}}}
\newunicodechar{ˣ}{\ensuremath{{}^{x}}}
\newunicodechar{ᵀ}{\ensuremath{{}^{\mathsf{T}}}}
\newunicodechar{ᵇ}{\ensuremath{{}^{b}}}
\newunicodechar{⁻}{\ensuremath{{}^{-}}}
\newunicodechar{⁺}{\ensuremath{{}^{+}}}
\newunicodechar{₀}{\ensuremath{{}_{0}}}
\newunicodechar{₁}{\ensuremath{{}_{1}}}
\newunicodechar{₂}{\ensuremath{{}_{2}}}
\newunicodechar{₃}{\ensuremath{{}_{3}}}
\newunicodechar{ᵢ}{\ensuremath{{}_{i}}}
\newunicodechar{ₐ}{\ensuremath{{}_{a}}}
% accented letters
\newunicodechar{î}{\ensuremath{\hat{\imath}}}
\newunicodechar{ĵ}{\ensuremath{\hat{\jmath}}}
\newunicodechar{ŷ}{\ensuremath{\hat{y}}}
\newunicodechar{Ŷ}{\ensuremath{\hat{Y}}}
\newunicodechar{ȳ}{\ensuremath{\bar{y}}}
% private-use slots written by prep.py for combining-accent sequences
\newunicodechar{}{\ensuremath{\hat{f}}}
\newunicodechar{}{\ensuremath{\hat{\beta}}}
\newunicodechar{}{\ensuremath{\hat{\sigma}}}
\newunicodechar{}{\ensuremath{\hat{k}}}
\newunicodechar{}{\ensuremath{\hat{\mu}}}
\newunicodechar{}{\ensuremath{\bar{x}}}

% --- layout -----------------------------------------------------------------
\setlength{\parskip}{0.45em}
\setlength{\parindent}{0pt}
\setlist{itemsep=0.15em,parsep=0.15em,topsep=0.3em}

\titleformat{\section}{\normalfont\Large\bfseries}{\thesection}{0.7em}{}
\titleformat{\subsection}{\normalfont\large\bfseries}{\thesubsection}{0.6em}{}
\titleformat{\subsubsection}{\normalfont\normalsize\bfseries}{\thesubsubsection}{0.5em}{}
\titlespacing*{\section}{0pt}{2.0ex plus .6ex minus .2ex}{1.0ex plus .2ex}
\titlespacing*{\subsection}{0pt}{1.6ex plus .5ex minus .2ex}{0.7ex plus .2ex}
\titlespacing*{\subsubsection}{0pt}{1.2ex plus .4ex minus .2ex}{0.5ex plus .2ex}

\pagestyle{fancy}
\fancyhf{}
\renewcommand{\headrulewidth}{0.4pt}
\fancyhead[L]{\footnotesize\nouppercase\leftmark}
\fancyhead[R]{\footnotesize\textit{\notestitle}}
\fancyfoot[C]{\footnotesize\thepage}
\fancypagestyle{plain}{\fancyhf{}\renewcommand{\headrulewidth}{0pt}\fancyfoot[C]{\footnotesize\thepage}}

% keep display math from being torn off its paragraph too eagerly
\allowdisplaybreaks
\setlength{\abovedisplayskip}{0.6em}
\setlength{\belowdisplayskip}{0.6em}
\setlength{\emergencystretch}{3em}
\sloppy
\ifLuaTeX
  \usepackage{selnolig}  % disable illegal ligatures
\fi
\IfFileExists{bookmark.sty}{\usepackage{bookmark}}{\usepackage{hyperref}}
\IfFileExists{xurl.sty}{\usepackage{xurl}}{} % add URL line breaks if available
\urlstyle{same}
\hypersetup{
  pdftitle={Problems ��� Implementation Notes},
  colorlinks=true,
  linkcolor={Maroon},
  filecolor={Maroon},
  citecolor={Blue},
  urlcolor={NavyBlue},
  pdfcreator={LaTeX via pandoc}}

\title{Problems ��� Implementation Notes}
\author{}
\date{}

\begin{document}
\maketitle

{
\hypersetup{linkcolor=black}
\setcounter{tocdepth}{2}
\tableofcontents
}
Idea, complexity, and notes for each named problem in the
\texttt{Problems/} package --- the recursion exercises, the DP problems,
grid shortest path, Held-Karp TSP, and the bipartite max-flow problems.
Reconciled against the code; \textbf{Code flag} lines mark
implementation deviations and known bugs (details in
\texttt{01\_Mismatch\_report.md}).

\emph{(The five recursion/divide-and-conquer exercise files ---
\texttt{Recursive\_multiplication}, \texttt{Recursive\_list\_sum},
\texttt{Recursive\_string\_reversal}, \texttt{Recursive\_max\_2d},
\texttt{Three\_way\_min} --- are documented here per the reorganization:
they are exercise-style leaf files, recommended to move from
\texttt{Algorithms/} to \texttt{Problems/}.)}

\hypertarget{recursion}{%
\subsection{Recursion}\label{recursion}}

A function that calls itself, used when a problem can be broken into
smaller subproblems of the same shape. Each call solves a smaller
instance until the problem is small enough to answer directly.

\begin{equation*}
T(n)=a\,T(n/b)+f(n)
\qquad\text{or}\qquad
T(n)=T(n-1)+f(n)
\end{equation*}

\begin{equation*}
\text{total work}=\sum_{\text{levels}}(\text{work per level});\quad
\text{stack depth}=\text{recursion depth}
\end{equation*}

\hypertarget{idea}{%
\subsubsection{Idea}\label{idea}}

Every recursive function has three pieces:

\begin{itemize}
\tightlist
\item
  \textbf{Base case} --- the terminating condition that stops the
  recursion. It can be implicit (e.g.~falling through when the input is
  empty). Without a reachable base case you get infinite recursion →
  stack overflow.
\item
  \textbf{Recursive call} --- the function calls itself on a
  \textbf{strictly smaller} subproblem, moving closer to the base case.
  This shrink is the crucial part: if the subproblem doesn't get smaller
  each time, the recursion never terminates.
\item
  \textbf{Body / work} --- the actual computation at each level:
  whatever this step contributes, combined with the result(s) of the
  recursive call(s).
\end{itemize}

\hypertarget{complexity}{%
\subsubsection{Complexity}\label{complexity}}

Time = (number of recursive calls) × (work per call); a recurrence like
\texttt{T(n)\ =\ a·T(n/b)\ +\ f(n)} captures it, and the Master Theorem
solves the common divide-and-conquer shapes. Space is O(maximum
recursion depth) for the call stack --- every pending call holds a frame
until it returns, so space can be significant even when total work is
small. Deep or unbounded recursion risks stack overflow; a
tail-recursive or iterative rewrite (or an explicit stack) reduces the
stack cost --- though Java has no tail-call optimization, so even a
tail-recursive form keeps its frames.

\hypertarget{applications}{%
\subsubsection{Applications}\label{applications}}

Divide and conquer (merge sort, quicksort, binary search), tree and
graph traversal (DFS), backtracking, memoized recursion / dynamic
programming, generating permutations and combinations, and problems
defined recursively by nature (factorial, Fibonacci, GCD). The four
\texttt{Recursive\_*} exercises below practice the pattern; merge sort
and DFS (see \texttt{Algorithms\_notes.md}) are where it graduates into
real algorithms.

\begin{center}\rule{0.5\linewidth}{0.5pt}\end{center}

\hypertarget{recursive_multiplication}{%
\subsection{Recursive\_multiplication}\label{recursive_multiplication}}

Multiply two integers using recursion, without the \texttt{*} operator.

\begin{equation*}
a\times b=\underbrace{a+a+\dots+a}_{b\text{ times}}
\;\Longrightarrow\;
T(b)=T(b-1)+O(1)=O(b)
\end{equation*}

\hypertarget{idea-1}{%
\subsubsection{Idea}\label{idea-1}}

\textbf{As implemented, the recursion is on \texttt{a} (the first
operand), and the workhorse is the doubling form:}

\begin{itemize}
\tightlist
\item
  Base case: \texttt{a\ ==\ 0\ →\ 0}.
\item
  Positive \texttt{a}:
  \texttt{multiply(a,\ b)\ =\ 2·multiply(a/2,\ b)\ +\ (a\ odd\ ?\ b\ :\ 0)}
  --- where the doubling itself is written as \texttt{m\ +\ m} (an
  addition, keeping the ``no \texttt{*}'' rule honest). Each call halves
  \texttt{a}, so the depth is \texttt{log₂(a)}.
\item
  Negative \texttt{a}: peel toward zero one step at a time ---
  \texttt{multiply(a,\ b)\ =\ multiply(a+1,\ b)\ −\ b}. Correct for any
  sign of \texttt{b}, but \textbf{linear}: depth is
  \texttt{\textbar{}a\textbar{}}.
\end{itemize}

\texttt{b} is never recursed on and never restricted --- it just rides
along and gets added/subtracted.

\hypertarget{complexity-1}{%
\subsubsection{Complexity}\label{complexity-1}}

Positive \texttt{a}: O(log a) time and stack depth. Negative \texttt{a}:
O(\textbar a\textbar) time and depth --- a large negative first operand
overflows the call stack. Returns \texttt{long} so \texttt{int\ ×\ int}
results don't wrap.

\hypertarget{notes}{%
\subsubsection{Notes}\label{notes}}

\begin{itemize}
\tightlist
\item
  \textbf{The asymmetry is the real lesson:} the doubling trick was
  applied to the positive branch only. The clean fix is to normalize the
  sign first (recurse on \texttt{\textbar{}a\textbar{}}, apply the sign
  at the top) so \emph{both} signs get the O(log) form --- or recurse on
  \texttt{min(\textbar{}a\textbar{},\ \textbar{}b\textbar{})} if you
  also want protection against a huge second operand in a linear form.
\item
  Doubling works because of binary decomposition:
  \texttt{a·b\ =\ (a/2)·b\ +\ (a/2)·b\ +\ (a\ odd\ ?\ b\ :\ 0)} is just
  reading \texttt{a}'s bits high-to-low.
\item
  The runner cross-checks against real \texttt{*} over a grid of signed
  values, which is what pinned down the negative-branch behavior.
\end{itemize}

\begin{center}\rule{0.5\linewidth}{0.5pt}\end{center}

\hypertarget{recursive_list_sum}{%
\subsection{Recursive\_list\_sum}\label{recursive_list_sum}}

Iterate an array with recursion instead of a loop --- sum all elements,
or sum only the odd-valued ones.

\hypertarget{idea-2}{%
\subsubsection{Idea}\label{idea-2}}

Recurse first, then contribute:
\texttt{f(arr,\ i)\ =\ f(arr,\ i+1)\ +\ g(arr{[}i{]})}, with base case
\texttt{i\ ==\ arr.length\ →\ 0} (an empty tail contributes nothing).
The index \texttt{i} is how the problem shrinks --- one step toward the
end per call; the additions happen on the unwind, so the sum accumulates
right-to-left (irrelevant for \texttt{+}, worth noticing for
non-commutative folds). A public method with no index wraps a private
helper that carries \texttt{i}, starting at 0.

The only thing that changes between ``sum all'' and ``sum odds'' is
\texttt{g}: for the total, \texttt{g(x)\ =\ x}; for odds, skip when
\texttt{x\ \%\ 2\ ==\ 0}, else add. Same traversal skeleton, different
per-element contribution.

\hypertarget{complexity-2}{%
\subsubsection{Complexity}\label{complexity-2}}

O(n) time and O(n) stack depth --- one frame per element until the base
case returns and the additions unwind. Sums accumulate in \texttt{long}.

\hypertarget{notes-1}{%
\subsubsection{Notes}\label{notes-1}}

\begin{itemize}
\tightlist
\item
  Skipping on \texttt{x\ \%\ 2\ ==\ 0} keeps negative odds
  (\texttt{−3\ \%\ 2\ ==\ −1\ ≠\ 0}) correctly included.
\item
  Because depth equals the array length, very large arrays overflow the
  stack --- this pattern is for learning the recursion shape, not
  production traversal. The two standard shapes are \textbf{head
  recursion} (recurse, then combine --- what this code does) and
  \textbf{tail recursion} (pass an accumulator down); Java keeps O(n)
  frames either way.
\end{itemize}

\begin{center}\rule{0.5\linewidth}{0.5pt}\end{center}

\hypertarget{recursive_string_reversal}{%
\subsection{Recursive\_string\_reversal}\label{recursive_string_reversal}}

Reverse a string using recursion.

\hypertarget{idea-3}{%
\subsubsection{Idea}\label{idea-3}}

\textbf{As implemented, each call peels \emph{both} ends:} take the last
character, the first character, and reverse the middle between them ---

\begin{verbatim}
reverse(s) = s[last] + reverse(s[1 .. last-1]) + s[0]
\end{verbatim}

Base case: a string of length 0 or 1 is its own reverse. Each call
shrinks the string by \textbf{two} characters, so the depth is ⌈n/2⌉ ---
half the frames of the classic peel-one-character forms
(\texttt{reverse(s{[}1:{]})\ +\ s{[}0{]}} or
\texttt{s{[}last{]}\ +\ reverse(s{[}:last{]})}), same idea.

\hypertarget{complexity-3}{%
\subsubsection{Complexity}\label{complexity-3}}

O(n) depth-wise? No --- \textbf{O(n²) time and memory}, because Java
strings are immutable: every \texttt{substring} and every \texttt{+}
copies a fresh string at each of the n/2 levels. Stack depth is ⌈n/2⌉.
To make it O(n): recurse over a \texttt{char{[}{]}} with two indices,
swap the ends, recurse inward (\texttt{reverse(arr,\ lo+1,\ hi−1)}, base
\texttt{lo\ \textgreater{}=\ hi}) --- O(1) work per level, no copying.
That two-pointer form is the array twin of exactly the both-ends peel
this code already does.

\hypertarget{notes-2}{%
\subsubsection{Notes}\label{notes-2}}

\begin{itemize}
\tightlist
\item
  The quadratic \texttt{substring}+\texttt{+} cost is the same
  immutability trap that makes string concatenation in a loop slow ---
  fine for learning, call it out.
\item
  Double-reversing any string returns the original --- a cheap invariant
  the runner tests.
\end{itemize}

\begin{center}\rule{0.5\linewidth}{0.5pt}\end{center}

\hypertarget{recursive_max_2d}{%
\subsection{Recursive\_max\_2d}\label{recursive_max_2d}}

Find the maximum value in a 2D array along with its \texttt{(row,\ col)}
position, recursively --- the ``bundle multiple returns into a result
object'' idea extended to a grid.

\hypertarget{idea-4}{%
\subsubsection{Idea}\label{idea-4}}

\textbf{As implemented:} recurse with an explicit \texttt{(i,\ j)} pair
that advances in row-major order --- bump the column; when it hits
\texttt{grid{[}0{]}.length}, wrap to column 0 of the next row. Each call
computes the best \texttt{Result} of the \textbf{suffix} (all cells
after \texttt{(i,\ j)}), then compares the current cell against it:

\begin{itemize}
\tightlist
\item
  Base case: \texttt{i} past the last row → return a sentinel
  \texttt{Result(Integer.MIN\_VALUE,\ −1,\ −1)} (an empty suffix can't
  win against any real cell).
\item
  Combine: keep the current cell when
  \texttt{suffixBest.value\ \textless{}=\ grid{[}i{]}{[}j{]}} --- the
  \texttt{\textless{}=} means the \textbf{current (earlier) cell wins
  ties}, so the earliest cell in row-major order is reported: lowest row
  first, then lowest column. That's the same answer a
  \texttt{\textgreater{}}-based nested-loop brute force produces, which
  is what the runner cross-checks.
\end{itemize}

This is one frame per cell --- neither the flattened single-index form
(\texttt{k\ →\ (k/C,\ k\%C)}) nor the two-level row/column recursion,
though it's closest in spirit to the flattened one with the decode
replaced by explicit carry logic.

\hypertarget{complexity-4}{%
\subsubsection{Complexity}\label{complexity-4}}

O(R·C) time --- every cell visited once. Stack depth is O(R·C) (one
frame per cell), so big grids will exhaust the stack; the two-level
row/column form would cut depth to O(R + C). One small \texttt{Result}
allocation per level.

\hypertarget{notes-3}{%
\subsubsection{Notes}\label{notes-3}}

\begin{itemize}
\tightlist
\item
  \textbf{Rectangular grids only:} the column wrap uses
  \texttt{grid{[}0{]}.length}, so a jagged array with rows shorter than
  row 0 indexes out of bounds (and longer rows are only partially
  scanned). The validation rejects empty grids and null rows but not
  jaggedness --- a two-level recursion would handle jagged rows
  naturally.
\item
  \textbf{Fix the tie-break policy explicitly} (first vs last, in which
  scan order) and test it --- here it's earliest-in-row-major, enforced
  by that single \texttt{\textless{}=}.
\end{itemize}

\begin{center}\rule{0.5\linewidth}{0.5pt}\end{center}

\hypertarget{three_way_min}{%
\subsection{Three\_way\_min}\label{three_way_min}}

Find the minimum of an array with a \textbf{three-way}
divide-and-conquer split --- the k-way-split exercise from the
divide-and-conquer primer (see \texttt{Algorithms\_notes.md}).

\hypertarget{idea-5}{%
\subsubsection{Idea}\label{idea-5}}

Split \texttt{{[}l,\ r{]}} into three roughly equal segments at
\texttt{m1\ =\ l\ +\ len/3} and \texttt{m2\ =\ l\ +\ 2·len/3}, recurse
into each, and combine by returning the smallest of the three
sub-minimums (a chain of \texttt{≤} comparisons that also happens to
favor the leftmost segment on ties --- irrelevant for a value-only min,
but the kind of thing that matters the moment you also return an index).

Base cases: one element returns itself; two elements return the smaller
--- needed because a 3-way split of a 2-element range would produce an
empty middle segment.

\hypertarget{complexity-5}{%
\subsubsection{Complexity}\label{complexity-5}}

\texttt{T(n)\ =\ 3T(n/3)\ +\ O(1)} → \textbf{O(n)} time (all n leaves
are visited --- no split count avoids that for min). Stack depth O(log₃
n).

\hypertarget{notes-4}{%
\subsubsection{Notes}\label{notes-4}}

\begin{itemize}
\tightlist
\item
  The point of the exercise is the recurrence shape, not speed --- a
  loop is obviously simpler and faster in constants.
\item
  \textbf{Code flag (C8):} the right recursion starts at
  \texttt{m2\ −\ 1}, so the middle and right segments overlap by one
  element. Harmless for min (overlap can't change a minimum, and no
  element is \emph{skipped}), but the split isn't a clean partition ---
  \texttt{min(m2,\ r)} is the intended boundary.
\item
  Null throws; an empty array is rejected with
  \texttt{IllegalArgumentException} --- min of nothing has no answer.
\end{itemize}

\begin{center}\rule{0.5\linewidth}{0.5pt}\end{center}

\hypertarget{dynamic-programming}{%
\subsection{Dynamic programming}\label{dynamic-programming}}

A technique for problems that break into \textbf{overlapping}
subproblems: solve each distinct subproblem once, store its answer, and
reuse it instead of recomputing. It applies when a problem has
\emph{optimal substructure} (the answer is built from answers to smaller
instances) and \emph{overlapping subproblems} (the same smaller
instances recur many times). That overlap is what separates DP from
divide-and-conquer --- D\&C subproblems are independent, so caching buys
nothing; DP subproblems repeat, so caching is the whole point.

\begin{equation*}
\text{runtime}\;=\;\underbrace{\#\text{states}}_{\text{size of the table}}\times\underbrace{\text{work per state}}_{\text{transitions}}
\end{equation*}

\hypertarget{idea-6}{%
\subsubsection{Idea}\label{idea-6}}

Four things to pin down: - \textbf{State} --- the parameters that
uniquely identify a subproblem (e.g.~``index \texttt{i}, remaining
capacity \texttt{w}''). - \textbf{Transition / recurrence} --- how a
state's answer is built from smaller states. - \textbf{Base cases} ---
the smallest states, answered directly. - \textbf{Evaluation order} ---
either \textbf{top-down} (write the natural recursion, add a memo table
so each state computes once --- how Mountain\_scenes,
Narrow\_art\_gallery, and Domino\_tromino\_tiling are written here) or
\textbf{bottom-up} (tabulation in dependency order --- how the tilings,
knapsacks, Magical\_cows, and TSP are written here).

\hypertarget{complexity-6}{%
\subsubsection{Complexity}\label{complexity-6}}

Roughly (number of distinct states) × (work per transition). Memory is
the table size, but it can often shrink: if a state only depends on the
previous row/day, keep just that (a \textbf{rolling array}) and drop the
full table.

\hypertarget{applications-1}{%
\subsubsection{Applications}\label{applications-1}}

Fibonacci, 0/1 knapsack, longest common subsequence, edit distance, coin
change, matrix-chain multiplication, and DP-flavored graph algorithms
(Bellman-Ford, Floyd-Warshall, Held-Karp TSP).

\begin{center}\rule{0.5\linewidth}{0.5pt}\end{center}

\hypertarget{magical_cows}{%
\subsection{Magical\_cows}\label{magical_cows}}

\textbf{Problem.} Farms hold cows, up to a capacity \texttt{C}. Every
night a farm with \texttt{v} cows doubles to \texttt{2v}; if
\texttt{2v\ \textless{}=\ C} it becomes one farm of \texttt{2v},
otherwise it splits into two farms of \texttt{v} cows each. Given the
initial farms and a list of query days, report the total number of farms
after each queried number of nights.

\hypertarget{idea-7}{%
\subsubsection{Idea}\label{idea-7}}

Simulating individual farms explodes --- farms can double every night,
so after \texttt{D} nights you may have up to \texttt{N·2\^{}D} of them.
The unlock: \textbf{individual farms don't matter, only how many cows
each holds}, and cow counts only ever range over \texttt{1..C}. So track
\texttt{cnt{[}v{]}} = number of farms currently holding \texttt{v} cows,
and advance the whole distribution one night at a time:

\begin{itemize}
\tightlist
\item
  for each \texttt{v} with \texttt{2v\ \textless{}=\ C}: those farms
  grow --- \texttt{next{[}2v{]}\ +=\ cnt{[}v{]}} (one farm each, now at
  \texttt{2v}).
\item
  for each \texttt{v} with \texttt{2v\ \textgreater{}\ C}: those farms
  split --- \texttt{next{[}v{]}\ +=\ 2\ *\ cnt{[}v{]}} (two farms each,
  still at \texttt{v}).
\end{itemize}

That's an O(C) step regardless of how many farms exist.

\textbf{As implemented:} the code builds the \textbf{full table}
\texttt{farm\_count{[}day{]}{[}v{]}} for every day from 0 up to the
largest query day (day 0 seeded from the initial farms, each next row
derived from the previous by the two rules above). A query is then
answered by \textbf{summing that day's row} --- O(C) per query, a plain
table lookup of the distribution rather than a precomputed total.

\hypertarget{complexity-7}{%
\subsubsection{Complexity}\label{complexity-7}}

O(C · D\_max) to build the table (D\_max = largest query day), then
\textbf{O(C) per query} (row sum). Space is \textbf{O(C · D\_max)} for
the full table. Both are easy to tighten if you ever care: keep only two
rolling rows (O(C) space) and record the running total per day as you go
(O(1) per query) --- the note-worthy optimizations, not what the code
currently does. Farm totals grow up to
\textasciitilde{}\texttt{N·2\^{}D}; counts are \texttt{long} (fits
comfortably for the usual constraints \texttt{D\ ≤\ 50},
\texttt{N,\ C\ ≤\ 1000}).

\hypertarget{notes-5}{%
\subsubsection{Notes}\label{notes-5}}

\begin{itemize}
\tightlist
\item
  The two transition cases are the crux: \texttt{2v\ \textless{}=\ C} →
  one farm at \texttt{2v}; \texttt{2v\ \textgreater{}\ C} → \textbf{two}
  farms at \texttt{v} (the split doubles that bucket's count).
\item
  Day 0 is just the initial distribution --- a query of 0 nights reads
  row 0 directly, which the table layout handles for free.
\item
  Validation: capacity ≥ 1, no negative query days, every initial herd
  in \texttt{{[}1,\ C{]}}.
\item
  Overflow is the lurking bug --- the table and sums are \texttt{long},
  not \texttt{int}.
\end{itemize}

\begin{center}\rule{0.5\linewidth}{0.5pt}\end{center}

\hypertarget{tiling}{%
\subsection{Tiling}\label{tiling}}

\textbf{Problem.} Count the number of ways to completely fill a 1×n
strip of slots using tiles of length 1 and length 2.

\begin{equation*}
f(n)=f(n-1)+f(n-2),
\qquad f(1)=1,\ f(2)=2
\qquad\text{(Fibonacci; }O(n)\text{ time, }O(1)\text{ space rolling)}
\end{equation*}

\hypertarget{idea-8}{%
\subsubsection{Idea}\label{idea-8}}

Look at how the \textbf{last slot} gets covered --- there are exactly
two disjoint, exhaustive possibilities: - a size-1 tile sits in the last
slot, leaving \texttt{n\ −\ 1} slots to fill → \texttt{ways(n\ −\ 1)}
ways, or - a size-2 tile covers the last two slots, leaving
\texttt{n\ −\ 2} slots → \texttt{ways(n\ −\ 2)} ways.

Since those cases don't overlap and cover everything,
\texttt{ways(n)\ =\ ways(n\ −\ 1)\ +\ ways(n\ −\ 2)}. Base cases:
\texttt{ways(0)\ =\ 1} (the empty strip has exactly one tiling --- place
nothing) and \texttt{ways(1)\ =\ 1}. That's the Fibonacci recurrence; in
fact \texttt{ways(n)\ =\ Fib(n\ +\ 1)}.

The subproblems overlap heavily --- plain recursion recomputes
\texttt{ways(n−2)}, \texttt{ways(n−3)}, \ldots{} exponentially many
times. \textbf{As implemented:} bottom-up tabulation in a
\texttt{long{[}n+3{]}} table (padded so tiny n never under-allocates),
seeded with \texttt{ways(0)=1,\ ways(1)=1,\ ways(2)=2}, filled left to
right.

\hypertarget{complexity-8}{%
\subsubsection{Complexity}\label{complexity-8}}

O(n) time, O(n) space for the table as written. Since each value depends
only on the previous two, two rolling variables would make it O(1) space
--- the table is kept for clarity. Values grow like Fibonacci (≈ φⁿ),
overflowing \texttt{long} around n = 91 --- \texttt{BigInteger} beyond
that.

\hypertarget{notes-6}{%
\subsubsection{Notes}\label{notes-6}}

\begin{itemize}
\tightlist
\item
  \textbf{\texttt{ways(0)\ =\ 1} is the load-bearing base case.} An
  empty strip has one tiling (do nothing); setting it to 0 makes every
  larger answer wrong --- the empty-product/empty-sum convention again.
\item
  \textbf{It's Fibonacci in disguise.} Recognizing that unlocks
  fast-doubling or matrix exponentiation for O(log n) if inputs get
  huge.
\item
  \textbf{Generalizes.} Tiles of sizes \texttt{\{1..k\}} give
  \texttt{ways(n)\ =\ ways(n−1)\ +\ …\ +\ ways(n−k)} --- which is
  exactly \texttt{Tiling\_general} with unit color counts.
\end{itemize}

\begin{center}\rule{0.5\linewidth}{0.5pt}\end{center}

\hypertarget{tiling_general}{%
\subsection{Tiling\_general}\label{tiling_general}}

\textbf{Problem.} Count the ways to fill a 1×n strip using tiles of
various lengths, where a given length may come in several colors. Tiles
arrive as parallel arrays: \texttt{lengths{[}i{]}} is an available
length, \texttt{colors{[}i{]}} how many distinct colors it comes in. A
tiling is a left-to-right sequence of colored tiles that exactly fills
n.

\hypertarget{idea-9}{%
\subsubsection{Idea}\label{idea-9}}

Same last-piece reasoning as basic tiling, generalized. Look at the
\textbf{last tile} placed: it has some length \texttt{L} and some color,
and it covers the final \texttt{L} slots, leaving \texttt{n\ −\ L}
before it. Sum over every way to choose that last tile:

\begin{verbatim}
ways(0) = 1
ways(m) = sum over i of  colors[i] * ways(m - lengths[i])   for lengths[i] <= m
\end{verbatim}

Two dimensions of choice fall out naturally: - \textbf{Different
lengths} → separate terms in the sum, each reaching back a different
distance. - \textbf{Different colors of the same length} → the
\texttt{colors{[}i{]}} multiplier: each color is a distinct way to place
that tile.

The basic 1-and-2 tiling is the special case
\texttt{lengths\ =\ \{1,\ 2\}}, \texttt{colors\ =\ \{1,\ 1\}}
(Fibonacci). A single length-1 tile with \texttt{k} colors gives
\texttt{k\^{}n}. Unit-color \texttt{\{1,\ 2,\ 3\}} gives Tribonacci.

\textbf{As implemented:} bottom-up over a \texttt{long{[}n+3{]}} table,
inner loop over tile types with the
\texttt{lengths{[}i{]}\ \textless{}=\ m} guard; inputs validated
(equal-length arrays, lengths and color counts ≥ 1, n ≥ 0).

\hypertarget{complexity-9}{%
\subsubsection{Complexity}\label{complexity-9}}

O(n · t) time, \texttt{t} = number of tile types. O(n) space for the
table (a rolling window of \texttt{maxLength} entries would suffice ---
the recurrence never reaches further back --- but the full table is
clearer).

\hypertarget{notes-7}{%
\subsubsection{Notes}\label{notes-7}}

\begin{itemize}
\tightlist
\item
  \textbf{\texttt{ways(0)\ =\ 1}} is still the anchor --- one way to
  tile nothing.
\item
  \textbf{The color count is a multiplier, not a new term.} Two colors
  of length 1 turn \texttt{+\ ways(m−1)} into \texttt{+\ 2·ways(m−1)}.
\item
  \textbf{Skip out-of-range tiles.} A tile longer than the remaining
  strip can't be the last piece --- the
  \texttt{lengths{[}i{]}\ \textless{}=\ m} guard.
\item
  \textbf{Watch overflow.} Up to \texttt{k\^{}n} growth; sums are
  \texttt{long}, \texttt{BigInteger} if the numbers can exceed it.
\item
  Representing colors as counts (rather than listing each colored tile
  separately) keeps the sum compact --- listing them would be
  arithmetically identical, just wasteful.
\end{itemize}

\begin{center}\rule{0.5\linewidth}{0.5pt}\end{center}

\hypertarget{domino_tiling_3xn}{%
\subsection{Domino\_tiling\_3xn}\label{domino_tiling_3xn}}

\textbf{Problem.} Count the ways to completely tile a 3×n board with 1×2
dominoes (horizontal or vertical).

\begin{equation*}
f(n)=4f(n-2)-f(n-4),
\qquad f(0)=1,\ f(2)=3
\qquad (f(n)=0 \text{ for odd } n)
\end{equation*}

\hypertarget{idea-10}{%
\subsubsection{Idea}\label{idea-10}}

\textbf{As implemented, this is a broken-profile (bitmask) column DP
with 8 states} --- the general technique, not the height-3 shortcut.

Fill the board column by column. The only thing the future needs to know
about the past is the \textbf{profile}: which of the 3 cells at the
current column boundary are already covered by dominoes sticking out of
the previous column. Three cells → a 3-bit mask → 8 states.
\texttt{dp{[}i{]}{[}s{]}} = number of ways to tile everything left of
column \texttt{i} such that column \texttt{i} has exactly the cells in
mask \texttt{s} pre-filled\ldots{} read at the answer end as:
\texttt{dp{[}i{]}{[}7{]}} = ways to tile the first \texttt{i} columns
flush (all 3 bits set = column fully covered, nothing protruding).

Transitions (hand-enumerated in the code, one line per legal way to
complete a column): - \textbf{Complement pairs} --- a cell left open in
the previous column's profile must be covered by a horizontal domino
poking into this column, so state \texttt{s} at column \texttt{i} draws
from state \texttt{7−s} at column \texttt{i−1}:
\texttt{0←7,\ 1←6,\ 2←5,\ 3←4,\ 4←3,\ 5←2,\ 6←1,\ 7←0}. -
\textbf{Vertical placements} --- within a column, a vertical domino can
cover two adjacent cells, adding the extra transitions: \texttt{3←7} and
\texttt{6←7} (a flush previous column plus one vertical pair in this
column, bottom-pair or top-pair) and \texttt{7←3,\ 7←6} (a protruding
pair completed by a vertical domino on the remaining two\ldots{}
i.e.~the mask gains two adjacent bits at once). The middle-plus-edge
combinations that would need a ``vertical'' over non-adjacent cells
simply have no transition --- that's how illegal placements are
excluded.

Base: \texttt{dp{[}0{]}{[}7{]}\ =\ 1} (zero columns, flush boundary ---
one way: do nothing). Answer: \texttt{dp{[}n{]}{[}7{]}}.

\hypertarget{complexity-10}{%
\subsubsection{Complexity}\label{complexity-10}}

O(n · 8) time and, as written, O(n · 8) space (a full table; two rolling
rows would be O(1)). Values grow
\textasciitilde{}\texttt{(2+√3)\^{}(n/2)} --- \texttt{long} overflows
eventually; \texttt{BigInteger} for large n.

\hypertarget{notes-8}{%
\subsubsection{Notes}\label{notes-8}}

\begin{itemize}
\tightlist
\item
  \textbf{Parity is automatic.} A 3×n board has \texttt{3n} cells, so
  odd n → 0; the mask DP produces those zeros on its own (no flush
  completion exists), no special-casing needed. The answer sequence runs
  \texttt{1,\ 0,\ 3,\ 0,\ 11,\ 0,\ 41,\ …}.
\item
  \textbf{Equivalent formulations} worth knowing: the two-state coupled
  form \texttt{f(n)\ =\ f(n−2)\ +\ 2g(n−1)},
  \texttt{g(n)\ =\ f(n−1)\ +\ g(n−2)} (f = flush, g = one-cell bump, the
  \texttt{2} = mirror-image bumps), and the even-n closed form
  \texttt{f(n)\ =\ 4f(n−2)\ −\ f(n−4)} with \texttt{f(0)=1,\ f(2)=3}.
  The 8-state mask version implemented here is the one that generalizes
  to any board height --- the profile just grows to \texttt{2\^{}height}
  states.
\item
  Independent check: exhaustive backtracking confirms small n --- kept
  in the runner, since hand-coded transition tables are easy to
  mis-transcribe.
\end{itemize}

\begin{center}\rule{0.5\linewidth}{0.5pt}\end{center}

\hypertarget{domino_tromino_tiling}{%
\subsection{Domino\_tromino\_tiling}\label{domino_tromino_tiling}}

\textbf{Problem.} Count the ways to fully tile a 2×n board using 2×1
dominoes and L-shaped trominoes (a 2×2 square minus one corner, any of 4
orientations). LeetCode 790; answers mod 1e9+7.

\hypertarget{idea-11}{%
\subsubsection{Idea}\label{idea-11}}

Trominoes let a column boundary be \emph{ragged} --- one cell of a
column filled while the other still waits on a piece from the next
column. \textbf{As implemented: top-down memoized recursion over the
state \texttt{(column\ i,\ top\ cell\ free?,\ bottom\ cell\ free?)}} ---
the profile DP directly, with a \texttt{dp{[}n{]}{[}4{]}} memo (−1 =
uncomputed) and the two booleans encoding the current column's
raggedness:

\begin{itemize}
\tightlist
\item
  Both cells free (state ``full choice''): place a vertical domino (→
  next column untouched), two horizontal dominoes (→ next column both
  cells taken\ldots{} expressed as recursing with the next column's
  flags), or one of two mirror-image trominoes (→ next column has one
  cell taken, top or bottom).
\item
  Exactly one cell free (two mirror states): finish with a horizontal
  domino into the next column, or a tromino that also bites one cell of
  the next column.
\item
  No cells free: this column is done --- advance with the next column
  fully free.
\end{itemize}

Base case: \texttt{i\ ==\ n} with a clean boundary → 1 tiling. Each
branch is guarded so nothing pokes past column n−1. Results are stored
mod 1e9+7; each state's sum is at most a handful of sub-results, so
\texttt{long} never overflows before the mod.

\hypertarget{complexity-11}{%
\subsubsection{Complexity}\label{complexity-11}}

O(n) states × O(1) transitions → O(n) time, O(n) memo space (the closed
form below runs in O(1) space, kept as a cross-check).

\hypertarget{notes-9}{%
\subsubsection{Notes}\label{notes-9}}

\begin{itemize}
\tightlist
\item
  \textbf{The closed recurrence \texttt{f(n)\ =\ 2f(n−1)\ +\ f(n−3)}}
  with \texttt{f(0)=1,\ f(1)=1,\ f(2)=2} (sequence
  \texttt{1,\ 1,\ 2,\ 5,\ 11,\ 24,\ 53,\ 117,\ …}) is what the state
  recursion collapses to --- useful for verifying, and what you'd
  hand-roll once you trust it. It is \emph{not} what the code runs; the
  state form is the one that transfers to wider boards and other tile
  sets.
\item
  \textbf{No parity shortcut here.} Unlike 3×n dominoes, a 2×n board
  with trominoes is tileable for every n ≥ 0.
\item
  \textbf{Mod discipline:} mod on store; all terms non-negative, so no
  wrap-around fix needed.
\item
  Independent check: exhaustive backtracking that lays each piece
  covering the first empty cell confirms small n --- the branch set and
  its guards are easy to mis-enumerate.
\end{itemize}

\begin{center}\rule{0.5\linewidth}{0.5pt}\end{center}

\hypertarget{mountain_scenes}{%
\subsection{Mountain\_scenes}\label{mountain_scenes}}

\textbf{Problem} (Kattis ``scenes'', NAIPC 2016). A ribbon of length
\texttt{n} is cut into \texttt{w} columns, each filled bottom-up to an
integer height in \texttt{{[}0,\ h{]}}. Total ribbon used (sum of
heights) must be ≤ \texttt{n} --- she needn't use it all. A
\textbf{mountain} scene must be \emph{uneven}: all-equal columns make a
``plain,'' not a mountain. Count distinct mountain scenes, mod 1e9+7.

\hypertarget{idea-12}{%
\subsubsection{Idea}\label{idea-12}}

Count everything, then subtract the non-mountains --- cleaner than
counting mountains directly.

\textbf{Total scenes} = height assignments \texttt{(c\_1,\ …,\ c\_w)},
each \texttt{c\_i\ ∈\ {[}0,\ h{]}}, with \texttt{sum\ ≤\ n}. \textbf{As
implemented: top-down memoized recursion}
\texttt{count(column,\ ribbonLeft)} --- at each column try every height
\texttt{0..h}, recursing with the ribbon reduced; base cases: negative
ribbon → 0, past the last column → 1. Memo table
\texttt{dp{[}w+1{]}{[}n+1{]}}, −1 sentinel, mod on every addition.

\textbf{Flat scenes} (the ones to remove): every column equal to some
height \texttt{k}, using \texttt{w·k} ribbon, valid while
\texttt{w·k\ ≤\ n} --- there are \texttt{min(h,\ ⌊n/w⌋)\ +\ 1} of them
(computed in the code as \texttt{min(w·h,\ n)/w\ +\ 1}, same value).

\textbf{Answer} = \texttt{(total\ −\ flats\ +\ MOD)\ mod\ 1e9+7} --- the
\texttt{+\ MOD} before the final mod keeps the subtraction non-negative
after wrapping.

\hypertarget{complexity-12}{%
\subsubsection{Complexity}\label{complexity-12}}

\textbf{O(w · n · h)} time as written --- \texttt{w·n} states, each
summing over \texttt{h+1} heights --- and \textbf{O(w · n)} memo space.
Two standard tightenings, noted but not implemented: cap the ribbon
dimension at \texttt{min(n,\ w·h)} (the frame can't hold more, so larger
\texttt{n} only bloats the table), and replace the inner height loop
with a prefix sum over the previous column's row (each column reads a
contiguous window), dropping the time to O(w · min(n, w·h)).

\hypertarget{notes-10}{%
\subsubsection{Notes}\label{notes-10}}

\begin{itemize}
\tightlist
\item
  \textbf{``Uneven'' means subtract \emph{all} flats, not just the empty
  scene} --- every constant-height configuration, including all-zero and
  completely-full. The samples pin this down:
  \texttt{25\ 5\ 5\ →\ 6\^{}5\ −\ 6\ =\ 7770}.
\item
  \textbf{\texttt{w\ =\ 1} is always 0} --- a single column is trivially
  uniform. Good sanity check.
\item
  Sample checks (all in the runner): \texttt{25\ 5\ 5\ →\ 7770},
  \texttt{15\ 5\ 5\ →\ 6050}, \texttt{10\ 10\ 1\ →\ 1022},
  \texttt{4\ 2\ 2\ →\ 6}.
\item
  Implementation notes: state lives in static fields (memo, W/H/N) ---
  fine solo, not reentrant; validation requires \texttt{w,\ h\ ≥\ 1},
  \texttt{n\ ≥\ 0}.
\end{itemize}

\begin{center}\rule{0.5\linewidth}{0.5pt}\end{center}

\hypertarget{narrow_art_gallery}{%
\subsection{Narrow\_art\_gallery}\label{narrow_art_gallery}}

\textbf{Problem} (Kattis ``narrowartgallery'', 2014 NAQ). A gallery has
N rows of 2 rooms (left, right), each with a value. Exactly \texttt{k}
rooms must be closed. To keep the gallery walkable top-to-bottom you may
\textbf{not} close two rooms in the same row, nor two
diagonally-touching rooms in adjacent rows (left in one row and right in
the next, or vice versa). Choose closures leaving the \textbf{maximum
total value open}.

\hypertarget{idea-13}{%
\subsubsection{Idea}\label{idea-13}}

\textbf{As implemented: minimize the value \emph{closed}, then
complement} --- \texttt{answer\ =\ totalValue\ −\ minLost}. (Maximizing
open value directly is the equivalent framing; the complement is what
this code --- like Fiset's --- does, and it makes the recursion's
contribution per step a single room value instead of a row sum.)

Top-down recursion \texttt{minLost(row,\ closuresLeft,\ prevSide)} where
\texttt{prevSide\ ∈\ \{left-closed,\ free,\ right-closed\}} --- the only
fact about the past the no-diagonal rule cares about. Per row, three
moves: - \textbf{Close nothing} → lose 0 here, \texttt{k} unchanged,
next side = free. Allowed after anything. - \textbf{Close left} → lose
\texttt{rooms{[}row{]}{[}0{]}}, spend a closure, next side = left.
Allowed only when prevSide is free or left (left-after-right is the
forbidden diagonal). - \textbf{Close right} → symmetric, allowed when
prevSide is free or right.

Base cases: \texttt{closuresLeft\ ==\ 0} → 0 (the remaining rows all
stay open --- this is what makes it \emph{exactly} k, since the count
can't go below zero); rows exhausted with closures still owed → an
infeasibility sentinel so that branch never wins. Memo
\texttt{dp{[}row{]}{[}k{]}{[}side{]}} with −1 sentinel.

\hypertarget{complexity-13}{%
\subsubsection{Complexity}\label{complexity-13}}

O(N · k · 3) states, O(1) work each → \textbf{O(N·k)} time and memo
space. Values are small ints.

\hypertarget{notes-11}{%
\subsubsection{Notes}\label{notes-11}}

\begin{itemize}
\tightlist
\item
  \textbf{``At most one per row'' + ``no diagonal'' is the whole
  constraint} --- both fall out of the three-way side state; you never
  look back more than one row.
\item
  \textbf{Feasibility:} any \texttt{0\ ≤\ k\ ≤\ N} is achievable (close
  the same side in k rows), so a valid answer always exists in range ---
  enforced by the \texttt{k\ ≤\ rooms.length} validation.
\item
  \textbf{Code flag (C7):} the infeasibility sentinel is the literal
  \textbf{10000}. A legitimate closed-sum above 10000 (e.g.~200 rows of
  \textasciitilde100-value rooms) would lose to the ``infeasible''
  branch and inflate the answer. Use a large sentinel
  (\texttt{Integer.MAX\_VALUE/2}) --- the fix is one constant.
\item
  Sample checks (verified against an independent subset brute force in
  the runner): \texttt{17}, \texttt{17}, \texttt{102}.
\end{itemize}

\begin{center}\rule{0.5\linewidth}{0.5pt}\end{center}

\hypertarget{knapsack_01}{%
\subsection{Knapsack\_01}\label{knapsack_01}}

\textbf{Problem.} Given \texttt{n} items, each with a weight and a
value, and a knapsack of capacity \texttt{W}, pick a subset with total
weight ≤ \texttt{W} and maximal total value. Each item is taken
\textbf{at most once} (0/1 --- no fractions, no duplicates).

\begin{equation*}
dp[i][w]=\max\bigl(dp[i-1][w],\;\;v_i+dp[i-1][w-w_i]\bigr)
\qquad O(nW)\text{ time — \emph{pseudo}-polynomial}
\end{equation*}

\hypertarget{idea-14}{%
\subsubsection{Idea}\label{idea-14}}

The decision for each item is binary: take it or leave it. Process items
one at a time and track, for every capacity budget, the best value
achievable so far.

\texttt{dp{[}i{]}{[}w{]}} = best value using the first \texttt{i} items
with budget \texttt{w}. For item \texttt{i} (1-indexed, so
\texttt{weights{[}i-1{]}}):

\begin{verbatim}
dp[0][w] = 0                                     (no items -> no value)
dp[i][w] = dp[i-1][w]                            if weights[i-1] > w   (can't fit -> must skip)
         = max( dp[i-1][w],                      skip item i
                values[i-1] + dp[i-1][w - weights[i-1]] )   take item i
\end{verbatim}

Answer: \texttt{dp{[}n{]}{[}W{]}}. \textbf{As implemented:} the full 2D
table, filled row by row with \texttt{w} ascending --- safe with two
rows in play, because ``take'' always reads the \emph{previous} row.
Inputs are validated (equal-length arrays, non-negative
weights/values/capacity); zero items short-circuits to 0.

\hypertarget{complexity-14}{%
\subsubsection{Complexity}\label{complexity-14}}

O(n·W) time and O(n·W) space as written. The classic space squeeze ---
one rolling row, \textbf{capacity loop descending} --- gets O(W);
descending is what keeps each item single-use
(\texttt{dp{[}w\ −\ weight{]}} still means the previous item's row),
while ascending on one row would silently become \emph{unbounded}
knapsack. Worth internalizing even though this file keeps the 2D table
(its sibling needs the table anyway for reconstruction).

Note this is ``pseudo-polynomial'': \texttt{W} is a value, not an input
length, so the cost is exponential in \texttt{W}'s bit-length. Knapsack
is NP-hard; the DP is efficient only when \texttt{W} is modest.

\hypertarget{notes-12}{%
\subsubsection{Notes}\label{notes-12}}

\begin{itemize}
\tightlist
\item
  \textbf{The one-line 0/1-vs-unbounded switch} (rolling-row loop
  direction) is the interview classic riding on this problem.
\item
  \textbf{Fractional knapsack is a different problem} --- greedy by
  value/weight ratio is optimal there; for 0/1, greedy fails and you
  need the DP.
\item
  Sample checks (independently verified):
  \texttt{w\{1,3,4,5\}\ v\{1,4,5,7\}\ cap7\ →\ 9},
  \texttt{w\{2,3,4,5\}\ v\{3,4,5,6\}\ cap5\ →\ 7}.
\end{itemize}

\begin{center}\rule{0.5\linewidth}{0.5pt}\end{center}

\hypertarget{knapsack_01_items}{%
\subsection{Knapsack\_01\_items}\label{knapsack_01_items}}

\textbf{Problem.} Same as 0/1 knapsack, but return \textbf{both} the
maximum value and \emph{which items} achieve it.

\hypertarget{idea-15}{%
\subsubsection{Idea}\label{idea-15}}

Build the identical full 2D table, then walk it backward to recover the
choices:

\begin{verbatim}
w = W
for i from n down to 1:
    if dp[i][w] != dp[i-1][w]:      # value changed => item i-1 was taken
        record index i-1
        w -= weights[i-1]           # give back its weight
    # else: skipped, budget unchanged
\end{verbatim}

\texttt{dp{[}i{]}{[}w{]}} differs from \texttt{dp{[}i-1{]}{[}w{]}}
exactly when including item \texttt{i−1} beat excluding it --- the
optimum at this cell \emph{used} the item. \textbf{As implemented:}
indices are collected during the walk (descending) and then
\texttt{Arrays.sort}ed ascending --- same result as the
push-front/reverse approach. Value and items are returned together in a
small immutable \texttt{Result} (with \texttt{equals}/\texttt{hashCode},
so runners can compare directly).

\hypertarget{complexity-15}{%
\subsubsection{Complexity}\label{complexity-15}}

O(n·W) time and O(n·W) space --- the table is now mandatory (the O(W)
rolling row forgets the per-item history the walk needs). The backward
walk adds O(n), the sort O(n log n) --- noise.

\hypertarget{notes-13}{%
\subsubsection{Notes}\label{notes-13}}

\begin{itemize}
\tightlist
\item
  \textbf{The ``equal ⇒ skipped'' rule is a tie-break.} When taking and
  skipping tie, this convention skips --- producing one specific optimal
  set. The item set isn't unique in general, only the value is; so the
  runner validates the returned set (distinct indices, fits capacity,
  sums to the optimal value) rather than demanding one exact answer,
  except where the optimum is provably unique.
\item
  \textbf{Value-only vs items:} \texttt{Knapsack\_01.maxValue} and this
  method must agree on the value for every input --- a cheap consistency
  check the runner performs.
\end{itemize}

\begin{center}\rule{0.5\linewidth}{0.5pt}\end{center}

\hypertarget{grid_shortest_path}{%
\subsection{Grid\_shortest\_path}\label{grid_shortest_path}}

\textbf{Problem.} Given a rectangular grid where \texttt{0} is open and
\texttt{1} is a wall, find the fewest 4-directional moves from a start
cell to a target cell, stepping only on open cells. Return the move
count, or \texttt{-1} if unreachable. (Lives in \texttt{Problems/} ---
it's a named problem built on the BFS idea.)

\begin{equation*}
dp[r][c]=\text{cost}(r,c)+\min\bigl(dp[r-1][c],\;dp[r][c-1]\bigr)
\qquad O(RC)
\end{equation*}

\hypertarget{idea-16}{%
\subsubsection{Idea}\label{idea-16}}

The grid \emph{is} a graph, just implicit: each open cell is a vertex,
each pair of orthogonally adjacent open cells an edge. Every move costs
1, so this is unweighted shortest path → \textbf{BFS}, generating
neighbors on the fly from coordinates instead of building a graph.

\textbf{As implemented} --- the same mark-on-dequeue style as the graph
BFS (see \texttt{Breadth\_first\_search} in
\texttt{Algorithms\_notes.md}), with a \texttt{dist} grid instead of
parents:

\begin{verbatim}
if start or target is a wall: return -1
dist[*][*] = -1;  dist[start] = 0;  enqueue start
while queue not empty:
    (r, c) = dequeue                       # two parallel Queues carry r and c
    if not visited[r][c]:
        mark visited
        for each of the 4 directions:
            (nr, nc) = neighbor
            if in bounds, open:
                if dist[nr][nc] == -1: dist[nr][nc] = dist[r][c] + 1   # first sight wins
                enqueue (nr, nc)
return dist[tr][tc]        # still -1 if never reached
\end{verbatim}

\texttt{dist} is set at \textbf{first discovery} (guarded by
\texttt{==\ -1}), which preserves shortest distances for the same reason
the graph BFS's first-parent-wins does. The \texttt{-1} initialization
doubles as the unreachable sentinel, so the final read needs no special
case. There is no early exit on reaching the target --- the BFS runs to
completion and the answer is read out of the table (an
\texttt{if\ (r,c)==target\ return} on dequeue is a valid optimization,
just not taken).

\hypertarget{complexity-16}{%
\subsubsection{Complexity}\label{complexity-16}}

O(R·C) time and space --- each cell expands once, ≤ 4 neighbors each;
the duplicate-tolerant queue holds at most O(R·C) redundant entries
(constant-degree grid, so no blowup).

\hypertarget{notes-14}{%
\subsubsection{Notes}\label{notes-14}}

\begin{itemize}
\tightlist
\item
  \textbf{Reject wall endpoints up front} --- the code returns −1
  immediately if start or target is a wall (the BFS would conclude the
  same, the explicit check is clearer), after validating bounds and
  rectangularity.
\item
  \textbf{Bounds check every neighbor} before indexing --- edge cells
  are where the out-of-bounds access hides.
\item
  \textbf{Coordinates travel as two parallel queues} (\texttt{r\_index},
  \texttt{c\_index}) dequeued in lockstep --- the third option next to
  encoding \texttt{r*C\ +\ c} or a small cell object.
\item
  \textbf{Variants} slot in cleanly: 8-directional movement (extend the
  direction arrays), multiple sources (enqueue all at distance 0),
  weighted terrain (breaks the unit-cost assumption → Dijkstra / 0-1
  BFS). Path reconstruction would add a parent grid exactly like the
  graph BFS.
\end{itemize}

\begin{center}\rule{0.5\linewidth}{0.5pt}\end{center}

\hypertarget{traveling_salesman}{%
\subsection{Traveling\_salesman}\label{traveling_salesman}}

\textbf{Problem.} Given \texttt{n} cities and an \texttt{n\ ×\ n} cost
matrix, find the cheapest tour that starts at a city, visits every city
exactly once, and returns to the start (minimum-cost Hamiltonian cycle).
Works for asymmetric costs (\texttt{dist{[}i{]}{[}j{]}} may differ from
\texttt{dist{[}j{]}{[}i{]}}).

\begin{equation*}
dp[S][j]=\min_{i\in S\setminus\{j\}}\bigl(dp[S\setminus\{j\}][i]+w(i,j)\bigr)
\qquad O(2^{n}n^{2})\ \text{time},\ O(2^{n}n)\ \text{space}
\end{equation*}

\hypertarget{idea-17}{%
\subsubsection{Idea}\label{idea-17}}

Brute force tries all \texttt{(n−1)!} orderings --- hopeless past
\textasciitilde12 cities. \textbf{Held-Karp} memoizes over \emph{sets}
of visited cities: the cost of finishing depends only on which cities
are visited and where you stand, not the order behind you.

State: \texttt{memo{[}i{]}{[}mask{]}} = minimum cost of a path that
starts at city 0, has visited exactly the bitmask \texttt{mask} (bit
\texttt{i} set), and currently stands at \texttt{i}.

\textbf{As implemented --- bottom-up by subset size:}

\begin{verbatim}
base:  memo[i][ {0, i} ] = dist[0][i]   for each i != 0
for r = 3..n:                            # subset sizes
    for each mask with bitCount r containing city 0:
        for each next in mask (not 0):
            memo[next][mask] = min over end in mask\{0,next} of
                               memo[end][mask ^ (1<<next)] + dist[end][next]
answer: min over end != 0 of memo[end][fullMask] + dist[end][0]
\end{verbatim}

Iterating masks in increasing subset size guarantees every predecessor
state is finished before it's read. \texttt{notIN(i,\ mask)} is the
bit-test helper; \texttt{fullMask\ =\ (1\textless{}\textless{}n)\ −\ 1};
INF = \texttt{Integer.MAX\_VALUE/2} so \texttt{INF\ +\ dist} can't
overflow.

\hypertarget{complexity-17}{%
\subsubsection{Complexity}\label{complexity-17}}

O(2ⁿ · n²) time, O(2ⁿ · n) space --- practical to n ≈ 18--20. Still
exponential (TSP is NP-hard); just vastly better than factorial.

\hypertarget{notes-15}{%
\subsubsection{Notes}\label{notes-15}}

\begin{itemize}
\tightlist
\item
  \textbf{Fix the start at city 0.} A tour's cost is rotation-invariant;
  pinning the start removes the n-fold symmetry and anchors the base
  case. \texttt{n\ =\ 1} short-circuits to 0; \texttt{n\ =\ 2} falls out
  as \texttt{dist{[}0{]}{[}1{]}\ +\ dist{[}1{]}{[}0{]}}.
\item
  \textbf{Sentinel discipline:} halved MAX\_VALUE keeps
  \texttt{INF\ +\ dist} finite; unvisited memo cells default to 0 but
  are never read thanks to the subset-size ordering and the in-mask
  guards.
\item
  The table size is allocated via \texttt{(int)\ Math.pow(2,\ n)} ---
  works, but \texttt{1\ \textless{}\textless{}\ n} is the exact,
  idiomatic form.
\item
  \textbf{Code flag (B10):} after computing the cost, the method also
  runs a \textbf{tour reconstruction} loop --- whose comparison is
  inverted (it replaces the candidate when the incumbent is strictly
  \emph{better}, building a worst tour) --- and then \textbf{discards
  the result} (the \texttt{tour} array is never returned). Dead code
  hiding a bug no test can reach: either fix the comparison and return
  the tour, or delete the loop.
\item
  Sample checks (independently verified): symmetric 4-city →
  \texttt{80}, asymmetric 3-city → \texttt{3}.
\end{itemize}

\begin{center}\rule{0.5\linewidth}{0.5pt}\end{center}

\hypertarget{bipartite-matching-via-max-flow-notes}{%
\section{Bipartite Matching via Max Flow ---
Notes}\label{bipartite-matching-via-max-flow-notes}}

Two problems that look unrelated but share one shape: model them as a
bipartite graph, add a super-source and super-sink, and the maximum flow
\emph{is} the maximum matching. The pattern:
\texttt{source\ →\ left\ nodes\ (cap\ 1)\ →\ right\ nodes\ (edge\ where\ allowed)\ →\ sink\ (cap\ 1\ or\ a\ limit)}.

\hypertarget{mice_and_owls}{%
\subsection{Mice\_and\_owls}\label{mice_and_owls}}

\textbf{Problem.} Mice are scattered on a plane; owls will strike. Each
mouse can run at most a fixed distance and must reach a hole to survive.
Each hole has a capacity. Maximize mice saved (equivalently, minimize
eaten).

\begin{equation*}
\text{max matching} = \text{max flow with unit capacities}
\qquad
O(E\sqrt{V})\ \text{via Dinic's (Hopcroft--Karp bound)}
\end{equation*}

\hypertarget{idea-18}{%
\subsubsection{Idea}\label{idea-18}}

Bipartite max flow with capacities on the right side. \textbf{As
implemented:} one capacity matrix of size \texttt{mice\ +\ holes\ +\ 2},
with mice at indices \texttt{0..m−1}, holes at \texttt{m..m+h−1}, source
at \texttt{m+h}, sink at \texttt{m+h+1}: - \textbf{Source → each mouse},
capacity 1 (each mouse is one unit to route). - \textbf{Mouse → hole},
capacity 1, whenever
\texttt{distance(mouse,\ hole)\ \textless{}=\ maxDistance} (the
\texttt{canReach} predicate --- Euclidean, and \texttt{\textless{}=} so
a mouse exactly at its limit still makes it). - \textbf{Hole → sink},
capacity = the hole's capacity.

Max flow (computed with the repo's \texttt{Ford\_fulkerson}) = mice
simultaneously routable to holes without exceeding any hole =
\textbf{mice saved}. Mice eaten = total − flow.

\hypertarget{why-it-works}{%
\subsubsection{Why it works}\label{why-it-works}}

Each flow unit is one mouse traveling source → mouse → hole → sink. Unit
capacity on \texttt{source\ →\ mouse} uses each mouse once; the hole's
limit sits on its sink edge; integer capacities guarantee an integral
max flow, i.e.~an actual assignment. This is the classic reduction of
capacity-constrained matching to max flow.

\hypertarget{notes-16}{%
\subsubsection{Notes}\label{notes-16}}

\begin{itemize}
\tightlist
\item
  \textbf{Hole capacity lives on the hole → sink edge, not the mouse →
  hole edges} --- that's what caps a hole's intake no matter how many
  mice can reach it.
\item
  Comparing squared distances would avoid the \texttt{sqrt} for exact
  integer inputs; doubles are fine here.
\item
  Any max-flow routine works; the networks are small, so the DFS
  Ford-Fulkerson is plenty.
\end{itemize}

\hypertarget{elementary-math-variant-notes-only-no-implementation}{%
\subsection{Elementary Math (variant --- notes only, no
implementation)}\label{elementary-math-variant-notes-only-no-implementation}}

\textbf{Problem.} Given \texttt{N} expressions, each a pair
\texttt{(a,\ b)}, choose one operation per pair --- \texttt{a\ +\ b},
\texttt{a\ −\ b}, or \texttt{a\ ×\ b} --- so that all \texttt{N} results
are \textbf{distinct}. Decide feasibility (and, in the full problem,
output the choices).

\hypertarget{idea-19}{%
\subsubsection{Idea}\label{idea-19}}

This is \textbf{pure bipartite matching}, not a capacity problem: -
\textbf{Left nodes = the N expressions.} \textbf{Right nodes = the
candidate result values} --- each expression contributes at most 3, so ≤
3N distinct values; dedupe them. - \textbf{Edge expression → value} for
each result that expression can produce. - \textbf{Source → expression}
cap 1; \textbf{value → sink} cap 1 (each result value usable by at most
one expression).

A \textbf{perfect matching} --- max flow = N --- means every expression
gets a distinct result; flow \textless{} N means impossible. The matched
edges name the operations for constructive output.

\hypertarget{why-its-the-same-family-as-mice-and-owls}{%
\subsubsection{Why it's the same family as Mice and
Owls}\label{why-its-the-same-family-as-mice-and-owls}}

Both are ``assign each left item to a right slot, each slot used a
limited number of times, respecting an allowed-pairs relation.''
Mice/holes has capacities \textgreater{} 1 on the sink edges; Elementary
Math is capacity 1 everywhere (a pure matching). Same
source→left→right→sink skeleton; only the meaning of the right nodes and
the sink capacities differ.

\hypertarget{notes-17}{%
\subsubsection{Notes}\label{notes-17}}

\begin{itemize}
\tightlist
\item
  \textbf{Distinctness = unit capacity on the value → sink edge.} That
  single-use constraint \emph{is} ``all results distinct.''
\item
  The value set is small (≤ 3N): enumerate each expression's three
  results, dedupe, map each distinct value to an index. Watch
  \texttt{a\ −\ b} signs and \texttt{a\ ×\ b} magnitudes when hashing
  --- negative and large products are legitimately distinct values (use
  \texttt{long}).
\end{itemize}

\end{document}
